% ===== 原文 pp.6--7 =====
% tr-01 片 b：§2.1.3 在一般曲面上（后半）
% 编号公式 tags: 2.10--2.18

% [原文 p.6]
% 注：原文 p.5 末句「…and can be denoted as」的公式续写 $\psi_A(U)=Z_A(U)$ 印在 p.6 图 2 之后，
% 译文已把该句整句排在片 a 末尾，本片从图 2 与下一句开始。
\begin{center}
\begin{minipage}{\textwidth}
\centering
\includegraphics[width=0.55\textwidth]{fig2.png}\par\vspace{2pt}
\figtitle{2}{一个圆盘和一个三孔球面，边界上的和乐已指定}
\end{minipage}
\end{center}

因此，在下文中，我们把「配分函数 $Z$」与「波函数 $\psi$」这两种说法交替使用。这种用圆盘上的
路径积分（更一般地，用高维的球）来定义波函数的方法，是由 Hartle 与 Hawking 开创的；因此它
常被称为 Hartle--Hawking 波函数。

\begin{center}
\begin{minipage}{\textwidth}
\centering
\includegraphics[width=0.55\textwidth]{fig3.png}\par\vspace{2pt}
\figtitle{3}{圆盘的面积可以通过粘接一个圆柱面来改变}
\end{minipage}
\end{center}

那么，与圆盘相联系的波函数 $\psi_A(U)$ 是什么呢？首先注意，我们可以把面积为 $A$ 的圆柱面
粘到面积为 $A'$ 的圆盘上，从而得到一个面积为 $A+A'$ 的圆盘，见图 3。这意味着
\begin{equation}
  \psi_{A+A'}(U)=e^{-Ac_2}\psi_{A'}(U)
  \tag{2.10}
\end{equation}
其中 $c_2$ 是按 $c_2(R)$ 作用在 $\chi_R(U)$ 上的算符。因此只需确定 $\psi_{A=0}(U)$。当面积
为零时，$U$ 被迫成为单位元，因此 $\psi_{A=0}(U)=\alpha\delta(U)$，其中 $\delta(U)$ 是群 $G$
的群流形上单位元处的 $\delta$ 函数，$\alpha$ 是比例常数。我们可以写出
\begin{equation}
  \delta(U)=\sum_R d_R\chi_R(U)
  \tag{2.11}
\end{equation}
其中 $d_R$ 可以由正交归一性 (2.6) 求出：
\begin{equation}
  d_R=\int_G\delta(U)\chi_R(U)\,dU=\tr_R 1=\dim R .
  \tag{2.12}
\end{equation}
这样我们得到
\begin{equation}
  \psi_A(U)=\alpha\sum_R e^{-Ac_2(R)}(\dim R)\chi_R(U).
  \tag{2.13}
\end{equation}
% [Tutor Note: p.6, 式 (2.11)--(2.13)：这里省略了 Peter--Weyl 完备性关系。群流形上单位元处的
%   $\delta$ 函数按特征标展开正是 $\delta(U)=\sum_R d_R\chi_R(U)$（系数 $d_R$ 待定）；
%   两边乘 $\chi_{R'}(U)^*$ 并对 $U$ 积分，用 (2.6) 得
%   $d_{R'}=\int_G\delta(U)\chi_{R'}(U)^*dU=\chi_{R'}(e)^*=\dim R'$，这就是 (2.12)。
%   由 (2.10) 取 $A'=0$ 即 $\psi_A(U)=e^{-Ac_2}\psi_0(U)$，而 $e^{-Ac_2}$ 按 $R$ 逐项作用，
%   代入 (2.11) 便得 (2.13)。
%   易错点：$e^{-Ac_2(R)}$ 与 $\dim R$ 都在求和号内、逐 $R$ 取值；把 $\dim R$ 提到求和号外就错了。]

% [原文 p.7]
第二，此时把上面得到的圆柱面上的哈密顿量改写为边界和乐为 $U$、$V$ 的圆柱面上的振幅，是很有
用的：
\begin{equation}
  \psi_A(U,V)=\sum_R e^{-Ac_2(R)}\chi_R(U)\chi_R(V^{-1}).
  \tag{2.14}
\end{equation}

第三，注意任何二维曲面都可以切成若干块，使得每一块都是带三个洞的球面，见图 2。设面积为
$A$，三个洞周围的和乐为 $U$、$V$、$W$。波函数 $\psi_A(U,V,W)$ 是什么？关键性质是：当一个
圆盘被缝到一个洞上时，它就变成一个圆柱面，见图 4。用方程表示，这就成为
\begin{equation}
  \int_G\psi_A(U,V,W)\psi_{A'}(W^{-1})\,dW=\psi_A(U,V).
  \tag{2.15}
\end{equation}
% [Tutor Note: p.7, 式 (2.14)--(2.15)：把 (2.16) 与 (2.13) 代入 (2.15) 左端，可以用 $W$ 积分直接
%   验算右端的面积记号。$W$ 积分给出 $\int_G\chi_R(W)\chi_{R'}(W^{-1})dW=\delta_{RR'}$
%   （即 (2.6)，$W^{-1}$ 这个逆是必需的），于是左端
%   $=\sum_R e^{-(A+A')c_2(R)}\chi_R(U)\chi_R(V)$，即面积 $A+A'$ 的圆柱振幅。由此看：
%   (i) (2.15) 右端应记作 $\psi_{A+A'}(U,V)$，原文印作 $\psi_A$，漏掉了所缝圆盘的面积 $A'$；
%   (ii) (2.14) 的 $\chi_R(V^{-1})$ 是「传播子/热核」约定（第二个边界取相反定向），而
%   (2.16)(2.17) 用的是「所有边界同向」约定，两者相差一个逆。$\SU(2)$ 的特征标为实，
%   $\chi_R(V^{-1})=\chi_R(V)$，两种写法完全等价，故对本文结论无影响；若采用 (2.14)
%   的原样，则 (2.16) 中的 $\chi_R(V)$ 也须相应改成 $\chi_R(V^{-1})$ 才自洽。
%   结论：(ii) 属**约定不一致，不是笔误**；本段真正需要登记的笔误只有 (i)。
%   以上均已按 700 dpi 裁图复核印刷原文。正文照原文排印。]
\begin{center}
\begin{minipage}{\textwidth}
\centering
\includegraphics[width=0.45\textwidth]{fig4.png}\par\vspace{2pt}
\figtitle{4}{把圆盘缝到一个洞上，会使三孔球面变成一个圆柱面}
\end{minipage}
\end{center}

利用 (2.13) 与 (2.14) 可以求出，唯一的解是
\begin{equation}
  \psi_A(U,V,W)=\frac{1}{\alpha}\sum_R e^{-Ac_2(R)}(\dim R)^{-1}
  \chi_R(U)\chi_R(V)\chi_R(W).
  \tag{2.16}
\end{equation}

这一结果可以很容易地推广到面积为 $A$、亏格为 $g$、带 $n$ 个洞、和乐为 $U_{i=1,\dots,n}$ 的
任意曲面：只需把 $\psi_A(U,V,W)$ 粘接起来。答案是
\begin{equation}
  \psi_{A,g}(U_i)=\alpha^{2-2g-n}\sum_R e^{-Ac_2(R)}
  \frac{\prod_i\chi_R(U_i)}{(\dim R)^{2g-2+n}} .
  \tag{2.17}
\end{equation}
% [Tutor Note: p.7, 式 (2.17) 的三组特例是核对指数最快的办法：
%   $g=0,n=3$ 时 $\alpha^{2-0-3}=\alpha^{-1}$、$(\dim R)^{2g-2+n}=\dim R$，逐项回到 (2.16)；
%   $g=0,n=2$ 时 $\alpha^0$、$(\dim R)^0$，得圆柱振幅 $\sum_Re^{-Ac_2(R)}\chi_R(U)\chi_R(V)$
%   （可用来对照 (2.14) 的那个逆，见上一条批注）；
%   $g=1,n=0$ 时 $\alpha^0$、$(\dim R)^0$、空乘积 $=1$，得 $\sum_Re^{-Ac_2(R)}$，正是 (2.9) 的
%   环面配分函数（取 $A=TL$）。三个特例都通过，说明 $2-2g-n$ 这个幂次没抄错。]
注意这一最终答案自动满足图 5 所示的两个三穿刺球面缝接的结合性。

最后，我们简要讨论上面得到的各种量对 $\alpha$ 的依赖。在作用量 (2.1) 中我们可以加入一个
局部项
\begin{equation}
  \delta S=\beta\int d^2x\sqrt{\mathsf{g}}\,R
  \tag{2.18}
\end{equation}
其中 $R$ 是度规 $\mathsf{g}$ 的 Riemann 曲率。在没有穿刺的曲面上，它积分为 $\beta(2-2g)$，
因而是拓扑的。
